\section{Specializing exactness of a finite free complex}

\begin{lemma}[A contracting identity for an exact vector-space complex]
\label{lem:kn-specialization-contracting-identity}
\lean{EnvelopeProperty.Vanishing.SpecializationFiniteComplex.exists_contraction}
\leanok
Let $A\xrightarrow{f}B\xrightarrow{g}C$ be exact vector spaces over a
field.  There are linear maps $a:B\to A$ and $b:C\to B$ such that
$fa+bg=1_B$.
\end{lemma}
\begin{proof}
Choose a right inverse $s:\operatorname{im}(g)\to B$ of the
surjection induced by $g$, and a retraction
$\pi:C\to\operatorname{im}(g)$ of the inclusion.  Put $b=s\pi$.
Then $gbg=g$, so $u=1-bg$ has image in $\ker(g)=\operatorname{im}(f)$.
Choose a right inverse of $A\to\operatorname{im}(f)$ and compose it
with $u$ viewed as taking values in $\operatorname{im}(f)$ to define
$a$.  It satisfies $fa=u$, giving $fa+bg=1$.
\end{proof}

\begin{lemma}[A surjective contracting endomorphism detects exactness]
\label{lem:kn-specialization-surjective-contraction}
\lean{EnvelopeProperty.Vanishing.SpecializationFiniteComplex.exact_of_surjective_contraction}
\leanok
Let $A\xrightarrow{f}B\xrightarrow{g}C$ be vector spaces with
$gf=0$ and $B$ finite dimensional.  If $fa+bg$ is surjective for some
linear maps $a:B\to A$ and $b:C\to B$, then the complex is exact
at $B$.
\end{lemma}
\begin{proof}
Write $h=fa+bg$ and $Q=B/\operatorname{im}(f)$.  Since $gf=0$,
$h$ preserves $\operatorname{im}(f)$ and induces an endomorphism
$\bar h$ of $Q$.  Surjectivity of $h$ implies surjectivity of
$\bar h$.  As $Q$ is finite dimensional, $\bar h$ is injective.
If $gx=0$, then $hx=fax$, so $\bar h([x])=0$ and $[x]=0$.
Thus $x\in\operatorname{im}(f)$.  The reverse inclusion is $gf=0$.
\end{proof}

\begin{theorem}[A principal neighborhood of exact finite free fibres]
\label{thm:kn-specialization-finite-free-exactness}
\lean{EnvelopeProperty.Vanishing.SpecializationFiniteComplex.exists_field_exactness_neighborhood}
\leanok
Let $R$ be a commutative ring and let
$R^a\xrightarrow{F}R^b\xrightarrow{G}R^c$ be a complex of finite
free modules.  Suppose $R\to k$ is a surjective map to a field and
the induced complex over $k$ is exact at its middle term.  There is
$r\in R$ whose image in $k$ is nonzero such that, for every map
$R\to K$ to a field with $r\ne0$ in $K$, the induced original
matrix complex over $K$ is exact at its middle term.
\uses{lem:kn-specialization-contracting-identity,lem:kn-specialization-surjective-contraction}
\end{theorem}
\begin{proof}
Choose contraction matrices over $k$ by the first lemma, and lift
their entries to matrices $A,B$ over $R$ using surjectivity of
$R\to k$.  Put $H=FA+BG$ and $r=\det(H)$.  The reduction of $H$
is the identity, so $r$ reduces to $1$.  Over any field where $r$
is nonzero, the image of $H$ is invertible and hence surjective.
The original identity $GF=0$ is retained by the coefficient map.
The second lemma now gives the exactness of the induced complex.
\end{proof}

\begin{corollary}[One neighborhood for finitely many complex degrees]
\label{cor:kn-specialization-uniform-finite-free-exactness}
\lean{EnvelopeProperty.Vanishing.SpecializationFiniteComplex.exists_uniform_field_exactness_neighborhood}
\leanok
\uses{thm:kn-specialization-finite-free-exactness}
For a finite family of original three-term finite free complexes whose
residue fibres are exact over the same field-valued point, one scalar
with residue $1$ defines a principal open on which every field fibre
of every complex is exact.
\end{corollary}
\begin{proof}
Take the product of the finitely many scalars supplied by the theorem.
Its residue is the product of finitely many ones.  If its image in a
field is nonzero, the image of each factor is nonzero.  Hence the
theorem applies to each complex separately.  For an empty family the
product is $1$ and the assertion is vacuous.
\end{proof}

\begin{corollary}[Generic exactness from residue exactness]
\label{cor:kn-specialization-finite-free-generic-exactness}
\lean{EnvelopeProperty.Vanishing.SpecializationFiniteComplex.generic_exact_of_residue_exact}
\leanok
\uses{thm:kn-specialization-finite-free-exactness}
If $R$ is a domain with fraction field $K$, then the original matrix
complex is exact over $K$ whenever its fibre over a surjective
field-valued point is exact.
\end{corollary}
\begin{proof}
The scalar constructed by the theorem has residue $1$, so it is
nonzero in $R$.  The canonical coefficient map to the fraction field
is injective, so that scalar remains nonzero in $K$.  Apply the
theorem with this canonical map.
\end{proof}

\begin{remark}[The proper derived-section model remains a producer]
\label{rmk:kn-specialization-perfect-complex-frontier}
These are exactness theorems for the actual finite free matrix complex
and its actual coefficient maps.  To specialize the cohomology of a
proper family, one must construct an appropriate finite free complex
for its derived sections and verify the fibre comparison.  No such
model or comparison is an input masquerading as a theorem here.
\end{remark}

\begin{corollary}[Native homology vanishes on the same principal neighborhood]
\label{cor:kn-specialization-native-homology}
\lean{EnvelopeProperty.Vanishing.SpecializationFiniteComplex.exists_homology_zero_neighborhood,EnvelopeProperty.Vanishing.SpecializationFiniteComplex.generic_homology_isZero_of_residue_homology_isZero}
\leanok
\uses{thm:kn-specialization-finite-free-exactness,cor:kn-specialization-finite-free-generic-exactness}
Form the actual three-term complex of modules from the original two
matrices and its canonical coefficient maps.  If its native homology
vanishes over the original closed field-valued point, then its native
homology vanishes at every field point in one principal neighborhood.
Over a domain it also vanishes after the canonical map to the fraction
field.
\end{corollary}
\begin{proof}
In the category of modules, a three-term complex has zero native
homology if and only if the range of its first map is the kernel of
its second.  This is precisely exactness of the original pair of
linear maps.  Apply the exactness theorem on the principal
neighborhood, or its fraction-field corollary, and use the reverse
implication of the same homology criterion.
\end{proof}

\begin{corollary}[The original closed point and the original generic point]
\label{cor:kn-specialization-original-closed-native-homology}
\lean{EnvelopeProperty.Vanishing.SpecializationFiniteComplex.exists_closed_fibre_homology_zero_neighborhood,EnvelopeProperty.Vanishing.SpecializationFiniteComplex.generic_homology_isZero_of_closed_fibre_homology_isZero}
\leanok
\uses{cor:kn-specialization-native-homology}
Let $\mathfrak p$ be an original maximal ideal of $R$.  If the
native homology of the original matrix complex vanishes after the
canonical coefficient map $R\to\kappa(\mathfrak p)$, it vanishes
on every original residue fibre in a principal open containing
$\mathfrak p$.  If $R$ is a domain, it also vanishes after the
canonical coefficient map to any fraction field of $R$.
\end{corollary}
\begin{proof}
For a maximal ideal, the canonical map to its residue field is
surjective.  Apply the preceding corollary with that exact map.
Its scalar has image $1$ in $\kappa(\mathfrak p)$, so it is
outside $\mathfrak p$.  For any other prime $\mathfrak q$ in
the resulting principal open, the image of the scalar in
$\kappa(\mathfrak q)$ is nonzero.  This gives the original
residue-fibre assertion; the generic assertion uses the canonical
map to the fraction field as in the preceding corollary.
\end{proof}
